In \cite{awodey2026},
a new approach to the semantics of identity types was introduced,
in a category with finite limits with a bipointed exponentiable object $I$,
playing the role of an interval. 
From a small number of easy-to-check assumptions on a universe,
including the presence of what we will call ``path types'',
this interval can be used to produce an elimination principle for identity types.

Though not all models with intensional identity types have path types,
many of the standard models of Martin-L\"of type theory have them.
These examples include the model in groupoids \cite{hofmann1995},
the discrete model in presheaves,
Voevodsky's model in simplicial sets \cite{kapulkin2021},
and various cubical models.

In addition to simplifying the construction of identity types in models of
type theory,
path types also present identity types in a style more consistent
with that of $\Sigma$ and $\Pi$-type formers in algebraic type theory.
In algebraic type theory, a universe has $\Sigma$-types and $\Pi$-types
precisely when there are pullback squares classifying certain
universal constructions by the universe of type families
(cf. \Cref{def:alg-pi-type-mltt,def:alg-sig-type-mltt}).
\[ \begin{tikzcd}
      {P_{\uu}{\Ulow}} & \Ulow \\
      {P_{\uu}{\U}} & \U
      \arrow["\lam", from=1-1, to=1-2]
      \arrow["{{P_{\uu}{\uu}}}"', from=1-1, to=2-1]
      \arrow["\lrcorner"{anchor=center, pos=0.125}, draw=none, from=1-1, to=2-2]
      \arrow["\uu", from=1-2, to=2-2]
      \arrow["\Pi"', from=2-1, to=2-2]
    \end{tikzcd} 
    \quad \quad 
    \begin{tikzcd}
      Q & \Ulow \\
      {P_\uu\U} & \U
      \arrow["\pair", from=1-1, to=1-2]
      \arrow["{\uu \pcomp \uu}"', from=1-1, to=2-1]
      \arrow["\lrcorner"{anchor=center, pos=0.125}, draw=none, from=1-1, to=2-2]
      \arrow["\uu", from=1-2, to=2-2]
      \arrow["\Sigma"', from=2-1, to=2-2]
    \end{tikzcd}\]
A similar pullback square will be used to define path types.

A version of the argument presented here has been formalised in Lean,
as part of the HoTTLean project \cite{hua2025},
which is summarised in \Cref{sec:hottlean}.
% The assumptions for having Path types also produces
% a cubical filling structure on the universe,
% making any type family classified by the universe a cubical Kan fibration.
We present a variation of the proof in \cite{awodey2026} in a style
that is more consistent with the formalisation,
using $\pi$-preclans rather than a lex category,
and replacing the assumption of an exponentiable interval with
cylinder and path functors.
In the formalisation,
abstracting away some of the cartesian structure
from \cite{awodey2026} simplifies the formalisation in Lean.
This results in an almost syntactic presentation of the argument,
where the cylinder and path structures are controlled by a \emph{modality}
(see for example \cite{licata2016, licata2017, gratzer2020, shulman2023}).
% Unlike \cite{awodey2026},
% we do not assume the underlying category is lex.
As a result, the setup will be slightly more complicated,
but the abstractions should ultimately clarify some technical
details appearing in the argument.

\section{Cylinders and paths}

We fix a $\pi$-preclan $(\CC,\RR)$ (\Cref{def:pi-preclan}) with a terminal object
and a cylinder structure (\Cref{def:cylinder-structure}).
We can use the cylinder structure
to define an endofunctor on the slice $\CC / \Gamma$
for every object $\Gamma \in \CC$.
\[
   \cyl[\Gamma] : \CC / \Gamma \to \CC / \Gamma
\]
\[
  (d : X \to \Gamma) \mapsto (d \circ \pi : \cyl X \to \Gamma)
\]
This results in a {cylinder structure in each slice $\CC / \Gamma$}.

Let us denote the bicategory of large categories using $\tcat$.
For a category $\CC$,
let us write $S(-) : \CC \to \tcat$
for the \emph{covariant slice pseudofunctor} (in fact a strict 2-functor),
that takes an object $\Gamma \in \CC$ to the slice
$S(\Gamma) = \CC / \Gamma$ and a morphism $\sigma : \Delta \to \Gamma$
to the postcomposition functor
\[ S(\sigma) = \sigma_! : \CC / \Delta \to \CC / \Gamma\]
\[ (h : X \to \Delta) \mapsto (\sigma \circ h : X \to \Gamma)\]

\begin{proposition}
  Given a cylinder structure on a category $\CC$,
  the cylinders over each $\Gamma \in \CC$
  \[\cyl[\Gamma] : \CC / \Gamma \to \CC / \Gamma\]
  constitute a pseudotransformation
  \[ {\cyl} : S (-) \to S (-) \]
  with modifications
  \[\delta_0, \delta_1 : \id \to {\cyl}
  \quad \quad
  \pi : {\cyl} \to \id
  \quad \quad 
  \psi : \cyl \circ \cyl \to \cyl \circ \cyl\]
   satisfying the same equations as before,
   but now as modifications.
    % \[
    % \begin{tikzcd}
    %   \id & {\cyl} & \id & {{\cyl} \circ {\cyl}} & {{\cyl} \circ {\cyl}} \\
    %   & \id &&& {{\cyl} \circ {\cyl}}
    %   \arrow["{\delta_0}", from=1-1, to=1-2]
    %   \arrow[equals, from=1-1, to=2-2]
    %   \arrow["\pi"', from=1-2, to=2-2]
    %   \arrow["{\delta_0}"', from=1-3, to=1-2]
    %   \arrow[equals, from=1-3, to=2-2]
    %   \arrow["\psi", from=1-4, to=1-5]
    %   \arrow[equals, from=1-4, to=2-5]
    %   \arrow["\psi", from=1-5, to=2-5]
    % \end{tikzcd}
    % \]
    % \[ \begin{tikzcd}
    %   {\cyl} & {\cyl} & {\cyl} & {\cyl} & {{\cyl} \circ {\cyl}} & {{\cyl} \circ {\cyl}} \\
    %     {{\cyl} \circ {\cyl}} & {{\cyl} \circ {\cyl}} & {{\cyl} \circ {\cyl}} & {{\cyl} \circ {\cyl}} & {\cyl} & {\cyl}
    %     \arrow[equals, from=1-1, to=1-2]
    %     \arrow["{\delta_0 {\cyl}}"', from=1-1, to=2-1]
    %     \arrow["{{\cyl} \delta_0}", from=1-2, to=2-2]
    %     \arrow[equals, from=1-3, to=1-4]
    %     \arrow["{\delta_1 {\cyl}}"', from=1-3, to=2-3]
    %     \arrow["{{\cyl} \delta_1}", from=1-4, to=2-4]
    %     \arrow["\psi", from=1-5, to=1-6]
    %     \arrow["{\pi {\cyl}}"', from=1-5, to=2-5]
    %     \arrow["{{\cyl} \pi}", from=1-6, to=2-6]
    %     \arrow["\psi", from=2-1, to=2-2]
    %     \arrow["\psi", from=2-3, to=2-4]
    %     \arrow[equals, from=2-5, to=2-6]
    %   \end{tikzcd}\]
\end{proposition}

Pseudonaturality of ${\cyl} : S(-) \to S(-)$ says that for any
map $\sigma : \Delta \to \Gamma$ there is an isomorphism
\begin{equation}\label{eq:frobenius-identity}
\begin{tikzcd}
	{\CC / \Delta} & {\CC / \Gamma} \\
	{\CC / \Delta} & {\CC / \Gamma}
	\arrow["{\sigma_!}", from=1-1, to=1-2]
  \arrow["{\cyl[\Delta]}"', from=1-1, to=2-1]
  \arrow["{\cyl[\Gamma]}", from=1-2, to=2-2]
	\arrow["{\sigma_!}"', from=2-1, to=2-2]
\end{tikzcd} \end{equation}

\begin{example}\label{ex:exponentiable-bipointed}
Consider the assumptions of \cite{awodey2026},
where there is a bipointed exponentiable object
$\delta_0, \delta_1 : 1 \rightrightarrows I$.
We could instantiate a cylinder structure
by taking ${\cyl} X := X \times I$, 
where $I$ is the given interval,
and $X$ is an object in the category.
It follows that the cylinder over $\Gamma$ of some $X$ in the slice
$\CC / \Gamma$ is
${\cyl}_{\Gamma} X := X \times_\Gamma I_\Gamma$.
Then the pseudonaturality square looks like the topos-theoretic ``Frobenius identity'':
\[ \sigma_! (X \times_\Delta I_\Delta) \iso (\sigma_! X) \times_\Gamma I_\Gamma\]
\end{example}

% \begin{axiom}\label{ax:preclan}
%   There is a $\pi$-preclan (\Cref{def:pi-preclan}) $(\CC,\RR)$.
% \end{axiom}

Recall the definition of a partial right adjoint
from \Cref{def:partial-right-adjoint}.
\begin{definition}\label{def:path-functor}
  We say that a $\pi$-preclan $(\CC,\RR)$ with a terminal object
  and a cylinder structure (\Cref{def:cylinder-structure})
  \emph{has path functors} if, additionally,
  for each object $\Gamma$ in $\CC$,
  the cylinder functor $\cyl[\Gamma] : \CC / \Gamma \to \CC / \Gamma$ over $\Gamma$
  has a partial right adjoint $\arr[\Gamma] : \RR(\Gamma) \to \RR(\Gamma)$.
  \[\begin{tikzcd}
    {\CC / \Gamma} & {\RR(\Gamma)} \\
    {\CC / \Gamma} & {\RR(\Gamma)}
    \arrow[""{name=0, anchor=center, inner sep=0}, "{\cyl[\Gamma]}"', from=1-1, to=2-1]
    \arrow[hook', from=1-2, to=1-1]
    \arrow[""{name=1, anchor=center, inner sep=0}, "{\arr[\Gamma]}"', from=2-2, to=1-2]
    \arrow[hook', from=2-2, to=2-1]
    \arrow["{{\dashv_\partial}}"{description}, draw=none, from=0, to=1]
  \end{tikzcd}\]
\end{definition}
We will call $\arr[\Gamma] A$ the \emph{paths} in $A$ over $\Gamma$,
and $\arr[\Gamma]$ the \emph{path functor}.

Partial right adjoints have counits (but not units on all elements);
we will denote the counit at an object $A \in \RR(\Gamma)$ by
\[ \epsilon_A : \cyl[\Gamma] \arr[\Gamma]{A} \to A\]
It is the morphism in $\CC / \Gamma$
that is conjugate to the identity $\id : \arr[\Gamma]{A} \to \arr[\Gamma]{A}$.

The $\pi$-preclan $\RR$ induces a pseudofunctor
$\RR(-) : \CC^\op \to \tcat$,
where the contravariant action on morphisms is given by pullback.

\begin{proposition}
  Suppose $(\CC,\RR)$ has path functors.
  The functors $\arr[\Gamma]$ together form a pseudotransformation
  $\arr : \RR(-) \to \RR(-)$.
  This pseudotransformation is equipped with modifications 
  \[\src, \trg : \arr \to {\id}
  \quad \quad
  \pi : {\id} \to \arr
  \quad \quad 
  \sym : \arr \circ \arr \to \arr \circ \arr\]
  satisfying the following equations
    \[
      \begin{tikzcd}
        & \id &&& {{\arr} \circ {\arr}} \\
        \id & \arr & \id & {{\arr} \circ {\arr}} & {{\arr} \circ {\arr}}
        \arrow[equals, from=1-2, to=2-1]
        \arrow["\rfl"{description}, from=1-2, to=2-2]
        \arrow[equals, from=1-2, to=2-3]
        \arrow["\src", from=2-2, to=2-1]
        \arrow["\trg"', from=2-2, to=2-3]
        \arrow[equals, from=2-4, to=1-5]
        \arrow["\sym"', from=2-4, to=2-5]
        \arrow["\sym"', from=2-5, to=1-5]
      \end{tikzcd}
    \]
    \[ \begin{tikzcd}
	{{\arr} \circ {\arr}} & {{\arr} \circ {\arr}} & {{\arr} \circ {\arr}} & {{\arr} \circ {\arr}} & \arr & \arr \\
	\arr & \arr & \arr & \arr & {{\arr} \circ {\arr}} & {{\arr} \circ {\arr}}
	\arrow["\sym", from=1-1, to=1-2]
	\arrow["{{\src {\arr}}}"', from=1-1, to=2-1]
	\arrow["{{{\arr} \src}}", from=1-2, to=2-2]
	\arrow["\sym", from=1-3, to=1-4]
	\arrow["{{\trg {\arr}}}"', from=1-3, to=2-3]
	\arrow["{{{\arr} \trg}}", from=1-4, to=2-4]
	\arrow[equals, from=1-5, to=1-6]
	\arrow["{{\rfl {\arr}}}"', from=1-5, to=2-5]
	\arrow["{{{\arr} \rfl}}", from=1-6, to=2-6]
	\arrow[equals, from=2-1, to=2-2]
	\arrow[equals, from=2-3, to=2-4]
	\arrow["\sym"', from=2-5, to=2-6]
\end{tikzcd}\]
\end{proposition}
\begin{proof}
  The pseudonaturality condition follows a standard
  (partial) adjoint hom-set bijection argument,
  making use of the Frobenius identity.
  Suppose $\sigma : \Delta \to \Gamma$,
  then naturally in $X \in \CC / \Delta$ we have,
  \begin{align*}
    & \CC / \Delta (X, \sigma^* \arr[\Gamma] A) \\
     \iso \; & \CC / \Gamma (\sigma_! X, \arr[\Gamma] A) \\
     \iso \; & \CC / \Gamma (\cyl[\Gamma] \sigma_! X, A) \\
     \iso \; & \CC / \Gamma (\sigma_! \cyl[\Delta] X, A) \\
     \iso \; & \CC / \Delta (\cyl[\Delta] X, \sigma^* A) \\
     \iso \; & \CC / \Delta (X, \arr[\Delta] \sigma^* A) \\
  \end{align*}
  From the Yoneda lemma,
  it follows that $\sigma^* \arr[\Gamma] A \iso \arr[\Delta] \sigma^* A$.
  This isomorphism holds in $\RR(\Delta)$ since it is a full
  subcategory of the slice.

  To define $\src : \arr[\Gamma] A \to A$ (and similarly $\trg$)
  for some $A \in \RR(\Gamma)$,
  consider the composition 
  \[ 
  \begin{tikzcd}
    {\arr[\Gamma]A} & {\cyl[\Gamma] \arr[\Gamma]{A}} \\
    & A
    \arrow["{\delta_0 }", from=1-1, to=1-2]
    \arrow["\src"', from=1-1, to=2-2]
    \arrow["{\epsilon_A}", from=1-2, to=2-2]
  \end{tikzcd}
  \]
  Under the partial adjunction,
  the map $\pi : \cyl[\Gamma] A \to A$ has a conjugate
  $\rfl : A \to \arr[\Gamma] A$.
  We can define
  $\sym : \arr[\Gamma] \arr[\Gamma] A \to \arr[\Gamma] \arr[\Gamma] A$
  via the following hom-set bijection
  \begin{align*}
    & \CC / X (\arr \arr A, \arr \arr A) & \ni \sym \\
    \iso \; & \CC / X (\cyl \arr \arr A, \arr A) \\
    \iso \; & \CC / X (\cyl \cyl \arr \arr A, A) \\
    \iso \; & \CC / X (\cyl \cyl \arr \arr A, A) & (\text{via } \psi) \\
    \iso \; & \CC / X (\cyl \arr \arr A, \arr A) \\
    \iso \; & \CC / X (\arr \arr A, \arr \arr A) & \ni \id
  \end{align*}
  
  We leave it to the reader to check that $\src, \trg, \rfl$,
  and $\sym$ are modifications
  that satisfy the required equations.
  Here, being a modification amounts to these operations being natural in
  $A \in \RR(\Gamma)$
  and stable (up to isomorphism) under pullback along maps $\Delta \to \Gamma$.
\end{proof}

For any $\RR$-object $A$ in the slice over $\Gamma$,
we obtain a \emph{pathobject factorisation} of the diagonal
$\Delta : A \to A \times_\Gamma A$ using the path functor.
\[\begin{tikzcd}
	{A} & {\arr[\Gamma]{A}} \\
	& {{A} \times_\Gamma A}
	\arrow["\rfl", from=1-1, to=1-2]
	\arrow["{\Delta}"', from=1-1, to=2-2]
	\arrow["{(\src,\trg)}", from=1-2, to=2-2]
\end{tikzcd}\]
% We will refer to the map $(\src,\trg) : \arr[\Gamma] A \to A$
% as the pathobject of $A$.

% Let us call the morphisms in $\RR$ fibrations.

% \begin{axiom}\label{ax:path-types-universe}
%   We have an $\RR$-algebraic universe (\Cref{def:algebraic-universe})
%   $\uu : \Ulow \to \U$.
%   We call $\uu$ \emph{the universe of small fibrations}.
%   Denoting the set of pullbacks of $\uu$ by $\RR_\U$,
%   we call maps in $\RR_\U$ \emph{small fibrations}.
%   We also require that every small fibration $A : X \to \Gamma$ in $\RR_\U$
%   is equipped with a chosen \emph{classifying map}
%   $\ulcorner A \urcorner : \Gamma \to \U$,
%   so that
%   \[ \begin{tikzcd}
% 	X & \Ulow \\
% 	\Gamma & \U
% 	\arrow[from=1-1, to=1-2]
% 	\arrow["A"', from=1-1, to=2-1]
% 	\arrow["\lrcorner"{anchor=center, pos=0.125}, draw=none, from=1-1, to=2-2]
% 	\arrow["\uu", from=1-2, to=2-2]
% 	\arrow["{{\ulcorner A \urcorner}}"', from=2-1, to=2-2]
% \end{tikzcd} \]
% \end{axiom}

Let us now consider the universal pathobject factorisation
on a universe $\uu : \Ulow \to \U$ in $\RR(\U)$.
We will say that $\uu : \Ulow \to \U$ is an $\RR$-algebraic
universe if it is an $\RR$-map and $\U$ is an $\RR$-object
(cf. \Cref{def:algebraic-universe} where $\RR$ was a clan).

\begin{definition}\label{def:mltt-alg-path-types-2}
  Suppose $(\CC,\RR)$ has path functors
  and an $\RR$-algebraic universe
  $\uu : \Ulow \to \U$.
  Then we will say that $\uu$ admits algebraic $\Path$-types
  when we have maps
  $\Path : \Ulow \times_\U \Ulow \to \U$
  and $\path : \arr[\U] \Ulow \to \Ulow$
  forming a pullback square
  \[
  \begin{tikzcd}
    {{\arr[\U] \Ulow}} & \Ulow \\
    {\Ulow \times_\U \Ulow} & \U
    \arrow["\path", from=1-1, to=1-2]
    \arrow["{(\src,\trg)}"', from=1-1, to=2-1]
    \arrow[from=1-2, to=2-2]
    \arrow["\Path"', from=2-1, to=2-2]
    \arrow["\lrcorner"{anchor=center, pos=0.125}, draw=none, from=1-1, to=2-2]
  \end{tikzcd}\]
\end{definition}

For an $\RR$-algebraic universe $\uu : \Ulow \to \U$.
Let us write $\RR_\U$ for the class of pullbacks of $\uu$.
\[
  \RR_\U := \{ A : X \to \Gamma \st \begin{tikzcd}
  X & \Ulow \\
  \Gamma & \U
  \arrow["\exists", dashed, from=1-1, to=1-2]
  \arrow["A"', from=1-1, to=2-1]
  \arrow["\lrcorner"{anchor=center, pos=0.125}, draw=none, from=1-1, to=2-2]
  \arrow["\uu", from=1-2, to=2-2]
  \arrow["\exists"', dashed, from=2-1, to=2-2]
\end{tikzcd}\}
\]
Let us call the maps in $\RR_\U$ \emph{small fibrations},
and the maps in $\RR$ fibrations.

\begin{proposition}\label{prop:universal-pathobject}
  Suppose $(\CC,\RR)$ has path functors
  and an $\RR$-algebraic universe.
  The following are equivalent
  \begin{itemize}
    \item For any small fibration $A \to \Gamma$,
      the pathobject projection
      $(\src,\trg) : \arr[\Gamma] A \to A \times_\Gamma A$
      is a small fibration.
    \item The universal pathobject projection
      $(\src,\trg) : \arr[\Gamma] \Ulow \to \Ulow \times_\U \Ulow$
      is a small fibration.
  \end{itemize}
\end{proposition}
If the universe admits $\Path$-types,
then certainly the conditions of \Cref{prop:universal-pathobject} hold.
Conversely, if the universe $\uu$ has classifying maps
(\Cref{def:alg-universe-classifying-map})
then the latter condition provides $\Path$-types on $\uu$.

% \begin{definition}\label{def:AM-with-path-types}
%   Let us call the combination of Axioms
%   \labelcref{ax:cylinder,ax:path-functor,ax:path-types-universe,,ax:universal-pathobject}
%   \[(\CC,\RR,\cyl,\arr,\uu : \Ulow \to \U, \Path, \path)\]
%   an algebraic model of MLTT with (a universe that admits) $\Path$-types.
% \end{definition}

\section{Hurewicz structure}

Let us fix a $\pi$-preclan $(\CC,\RR)$ that has path functors.
Recall the definition of a Hurewicz structure on a morphism
from \Cref{def:hurewicz-defs}.

\begin{lemma}\label{lem:hurewicz-over-1}
  If $\uu : \Ulow \to \U$ is an $\RR$-algebraic universe 
  that admits a normal Hurewicz structure $(\uu, l)$,
  then any small fibration $f : A \to B$ (i.e., any pullback of $\uu$)
  inherits a normal Hurewicz structure.
\end{lemma}
\begin{proof}
  The right lifting structure of Hurewicz maps is
  preserved under pullback.
  Thus, the Hurewicz structure on the universe pulls back to 
  Hurewicz structure on all small fibrations.
\end{proof}

% The following is closely related to opfibrancy in $\Cat$
% (\Cref{def:split-opfibration}).
% \begin{definition}\label{def:hurewicz-defs} Let $\Gamma$ be an object of $\CC$ and
%   $f : A \to B$ be a morphism in $\CC / \Gamma$.
%   A \emph{Hurewicz structure $(f,l)$ over $\Gamma$} consists of
%   \begin{itemize} 
%     \item For every object $X$ in the slice $\CC / \Gamma$,
%     and every pair of maps $a : X \to A$
%     and $p : \cyl[\Gamma] X \to B$ such that
%     $f \circ a = p \circ \delta_0$,
%     a chosen diagonal filler $l_{(a,p)} : \cyl[\Gamma] X \to A$.
%     \[\begin{tikzcd}
%         X & A \\
%         {\cyl[\Gamma] X} & B
%         \arrow["a", from=1-1, to=1-2]
%         \arrow["\delta_0"', from=1-1, to=2-1]
%         \arrow["f", from=1-2, to=2-2]
%         \arrow["{l_{(a,p)}}"{description}, dashed, from=2-1, to=1-2]
%         \arrow["p"', from=2-1, to=2-2]
%       \end{tikzcd} \]
%     \item For any morphism $x : X' \to X$,
%     a coherence equation between lifts
%     $l_{(a \circ x,p \circ \cyl x)} : \cyl[\Gamma] X \to A$
%     and $l_{(a,p)} : \cyl[\Gamma] X \to A$.
%     \[ l_{(a,p)} \circ \cyl x = l_{(a \circ x,p \circ \cyl x)} \]
%   \end{itemize}

%   We say that a Hurewicz structure
%   $(f,l)$ is \emph{normal} if
%   for every $b : X \to B$
%   \[\begin{tikzcd}
%       X && A \\
%       {\cyl[\Gamma] X} & X & B
%       \arrow["a", from=1-1, to=1-3]
%       \arrow["{\delta_0}"', from=1-1, to=2-1]
%       \arrow["f", from=1-3, to=2-3]
%       \arrow["{l_{(a,b \circ \pi)}}"{description}, dashed, from=2-1, to=1-3]
%       \arrow["\pi"', from=2-1, to=2-2]
%       \arrow["a"{description}, from=2-2, to=1-3]
%       \arrow["b"', from=2-2, to=2-3]
%     \end{tikzcd}
%     \quad \quad 
%     l_{(a,b \circ \pi)} = a \circ \pi
%     \]  
% \end{definition}
% When $\cyl = (-) \times I$ is given by product with an interval object,
% a Hurewicz structure on $f : A \to B$ is the usual homotopy lifting
% definition from \cite{awodey2026}.

\begin{proposition}[Locally Hurewicz]\label{prop:hurewicz-equiv-defs}
  Let $\Gamma$ be an object of $\CC$ and
  $f : A \to B$ be a morphism in $\RR(\Gamma)$.
  (Note that $A$ and $B$ are required to be fibrant over $\Gamma$.)
  The following are equivalent (structures)
  \begin{itemize}
    \item A Hurewicz structure $(f,l)$ in the slice $\CC / \Gamma$ with respect to
      the cylinder $\cyl[\Gamma] : \CC / \Gamma \to \CC / \Gamma$.
    \item A section $l$ of the ``pullback-hom'' $h$
      given in the following diagram in $\RR(\Gamma)$.
      \[\begin{tikzcd}
      {\arr[\Gamma] A} && \\
      & P & A \\
      & {\arr[\Gamma] B} & B
      \arrow["h"{description}, dashed, from=1-1, to=2-2]
      \arrow["\src", bend left, from=1-1, to=2-3]
      \arrow["{\arr[\Gamma] f}"', bend right, from=1-1, to=3-2]
      \arrow[from=2-2, to=2-3]
      \arrow[from=2-2, to=3-2]
      \arrow["\lrcorner"{anchor=center, pos=0.125}, draw=none, from=2-2, to=3-3]
      \arrow["f", from=2-3, to=3-3]
      \arrow["\src"', from=3-2, to=3-3]
      \end{tikzcd}
      \quad \quad
      \quad \quad
      \begin{tikzcd}
        & {\arr[\Gamma] A} \\
        P & P
        \arrow["h", from=1-2, to=2-2]
        \arrow["l", from=2-1, to=1-2]
        \arrow[equals, from=2-1, to=2-2]
      \end{tikzcd} \]
  \end{itemize}

  Then a Hurewicz structure $(f,l)$ is normal
  if and only if $l \circ h \circ \rfl = \rfl$.
  \[\begin{tikzcd}
    A & {\arr[\Gamma] A} \\
    {\arr[\Gamma] A} & P
    \arrow["\rfl", from=1-1, to=1-2]
    \arrow["\rfl"', from=1-1, to=2-1]
    \arrow["h"', from=2-1, to=2-2]
    \arrow["l"', from=2-2, to=1-2]
  \end{tikzcd}\]
\end{proposition}

\begin{lemma}[Globally Hurewicz implies locally Hurewicz]\label{lem:global-hurewicz-to-local-hurewicz}
  If a morphism $f : A \to B$ in $\RR(\Delta)$ is (normal) Hurewicz
  as a map in $\CC$,
  then $f$ is also (normal) Hurewicz as a map in the slice $\CC / \Delta$.
\end{lemma}
\begin{proof}
  Let us prove a more general statement:
  if $\sigma : \Delta \to \Gamma$ is a morphism in $\CC$
  and $f : A \to B$ in $\RR(\Delta)$
  is (normal) Hurewicz over $\Gamma$,
  then $f : A \to B$ is (normal) Hurewicz over $\Delta$.

  Consider a diagram in $\CC / \Delta$:
  \[\begin{tikzcd}
      X & A \\
      {\cyl[\Gamma] X} & B
      \arrow["a", from=1-1, to=1-2]
      \arrow["\delta_0"', from=1-1, to=2-1]
      \arrow["f", from=1-2, to=2-2]
      \arrow["p"', from=2-1, to=2-2]
    \end{tikzcd} \]
  We take the image of this diagram under the composition functor
  $\sigma_! : \CC / \Delta \to \CC / \Gamma$.
  By the Frobenius identity (\cref{eq:frobenius-identity}),
  we obtain a diagonal filler in $\CC / \Gamma$.
  \[
  \begin{tikzcd}
    {\sigma_!X} & {\sigma_!X} & {\sigma_!A} \\
    {\cyl[\Gamma]\sigma_! X} & {\sigma_!\cyl[\Gamma] X} & {\sigma_!B}
    \arrow[equals, from=1-1, to=1-2]
    \arrow["{\delta_0}"', from=1-1, to=2-1]
    \arrow["{\sigma_!a}", from=1-2, to=1-3]
    \arrow["{\sigma_!\delta_0}"{pos=0.7}, from=1-2, to=2-2]
    \arrow["{\sigma_!f}", from=1-3, to=2-3]
    \arrow["l"{pos=0.2}, dashed, from=2-1, to=1-3]
    \arrow["\iso"{description}, draw=none, from=2-1, to=2-2]
    \arrow["{\sigma_!p}"', from=2-2, to=2-3]
  \end{tikzcd}
  \]
  The composition of the isomorphism 
  ${\cyl[\Gamma]\sigma_! X} \iso {\sigma_!\cyl[\Gamma] X}$
  with $l$ provides us with the required diagonal filler over $\Delta$.
\end{proof}

\begin{corollary}\label{cor:small-fibration-hurewicz}
  If $\uu : \Ulow \to \U$ is an $\RR$-algebraic universe 
  with a normal Hurewicz structure $(\uu, l)$,
  and $A$ and $B$ are objects in $\RR(\Gamma)$
  then any small fibration $f : A \to B$
  is normal Hurewicz $(f,l)$ in the slice over $\Gamma$.
\end{corollary}
\begin{proof}
  This follows from \Cref{lem:hurewicz-over-1} and \Cref{lem:global-hurewicz-to-local-hurewicz}.
% we have that the ``global'' Hurewicz condition on the universe
% implies a ``local'' Hurewicz condition on each small fibration.
\end{proof}

In the following,
we will always work in the slice over $\Gamma \in \CC$,
so that we drop the subscript in $\arr[\Gamma] A$ and $A \times_\Gamma A$
for an object $A \in \RR(\Gamma)$.
\begin{lemma}[Connection]\label{lem:connection}
  Suppose $\uu : \Ulow \to \U$ is an $\RR$-algebraic universe 
  with a normal Hurewicz structure $(\uu, l)$,
  and $\Gamma$ is an object in $\CC$.
  For any small fibration $A \to \Gamma$,
  there is a map $\chi : \arr A \to \arr \arr A$
  satisfying
  \[ \begin{tikzcd}
        A & {\arr A} \\
        {\arr A} & {\arr \arr A} \\
        & {\arr A}
        \arrow["\rfl", from=1-1, to=1-2]
        \arrow["\src_A", from=2-1, to=1-1]
        \arrow["\chi", from=2-1, to=2-2]
        \arrow[equals, from=2-1, to=3-2]
        \arrow["\src_{\arr \! A}"', from=2-2, to=1-2]
        \arrow["\trg_{\arr \! A}", from=2-2, to=3-2]
    \end{tikzcd} 
    \quad \quad
    \text{and}
    \quad \quad
    \begin{tikzcd}
	A & {\arr A} \\
        {\arr A} & {\arr \arr A}
        \arrow["\rfl_A", from=1-1, to=1-2]
        \arrow["\rfl_A"', from=1-1, to=2-1]
        \arrow["\rfl_{\arr \! A}", from=1-2, to=2-2]
        \arrow["\chi"', from=2-1, to=2-2]
    \end{tikzcd}
    \]
\end{lemma}
We will call $\chi$ a (normal) connection,
where the square involing $\rfl$ is the normality condition.
We think of $\chi$ as taking a path $p : x \rightsquigarrow y$ in $A$
to a path $\chi_p : \rfl x \rightsquigarrow p$ in $\arr A$,
or a homotopy in $A$,
such that $\chi_{\rfl x}$ is the constant path $\rfl : \rfl x \rightsquigarrow \rfl x$:
\[ \begin{tikzcd}
	x & x \\
	x & y
	\arrow[rightsquigarrow, from=1-1, to=1-2]
	\arrow[""{name=0, anchor=center, inner sep=0}, "{\rfl x}"', rightsquigarrow, from=1-1, to=2-1]
	\arrow[""{name=1, anchor=center, inner sep=0}, "p", rightsquigarrow, from=1-2, to=2-2]
	\arrow[rightsquigarrow, from=2-1, to=2-2]
	% \arrow["\chi"{description}, between={0.1}{0.9}, Rightarrow, from=0, to=1]
  \end{tikzcd} \]
We will construct $\chi$ via a ``box-filling'' argument,
schematically as follows.
First, the open box
\[ \begin{tikzcd}
	x & x \\
	x & y
	\arrow["{\rfl x}"', rightsquigarrow, from=1-1, to=2-1]
	\arrow["{\rfl x}", rightsquigarrow, from=1-1, to=1-2]
	\arrow["p"', rightsquigarrow, from=2-1, to=2-2]
  \end{tikzcd} \]
is a Hurewicz lifting problem in $(\src,\trg) : \arr A \to A \times A$.
By \Cref{cor:small-fibration-hurewicz},
the map $\arr A \to A \times A$ is Hurewicz over $\Gamma$.
Hence, we have filled box
\[ \begin{tikzcd}
	x & x \\
	x & y
	\arrow["{\rfl x}"', rightsquigarrow, from=1-1, to=2-1]
	\arrow["{\rfl x}", rightsquigarrow, from=1-1, to=1-2]
	\arrow["p"', rightsquigarrow, from=2-1, to=2-2]
	\arrow["", rightsquigarrow, from=1-2, to=2-2]
  \end{tikzcd} \]
Using $\sym : \arr \arr A \to \arr \arr A$,
we obtain the required homotopy. 
\[ \begin{tikzcd}
	x & x \\
	x & y
	\arrow["{\rfl x}", rightsquigarrow, from=1-1, to=1-2]
	\arrow["{\rfl x}"', rightsquigarrow, from=1-1, to=2-1]
	\arrow["p", rightsquigarrow, from=1-2, to=2-2]
	\arrow[rightsquigarrow, from=2-1, to=2-2]
  \end{tikzcd} \]
Notice that although we will know that the top map is $\rfl x$
(definitionally),
we will not be able to say anything about the bottom map provided by box filling.
\begin{proof}[Proof of \Cref{lem:connection}.]
  The open box can be constructed as follows.
  Since $\arr : \RR(\Gamma) \to \RR(\Gamma)$ is a partial right adoint,
  it preserves all limits,
  and so $\arr (A \times A) \iso \arr A \times \arr A$.
  We construct the open box $b$
  a map into the pullback.
  \[ \begin{tikzcd}
        {\arr A} && A \\
        & P & {\arr A} \\
        {\arr A \times \arr A} & {\arr (A \times A)} & {A \times A}
        \arrow["\src", from=1-1, to=1-3]
        \arrow["b"{description}, dashed, from=1-1, to=2-2]
        \arrow["{(\rfl \circ \src, \id)}"', from=1-1, to=3-1]
        \arrow["\rfl", from=1-3, to=2-3]
        \arrow["\src", from=2-2, to=2-3]
        \arrow["{\arr (\src,\trg)}"', from=2-2, to=3-2]
        \arrow["{(\src,\trg)}", from=2-3, to=3-3]
        \arrow["\iso"{description}, draw=none, from=3-2, to=3-1]
        \arrow["\src"', from=3-2, to=3-3]
        \arrow["\lrcorner"{anchor=center, pos=0.125}, draw=none, from=2-2, to=3-3]
    \end{tikzcd} \]
  By \Cref{cor:small-fibration-hurewicz},
  the map $\arr A \to A \times A$ is Hurewicz over $\Gamma$.
  Then we can obtain a filled box by composing $b$ with $l : P \to \arr \arr A$
  (see \Cref{prop:hurewicz-equiv-defs}),
  and then flip the box by composing with $\sym$.
%   Then we can construct a map $b : \arr A \to P$
%   \[\begin{tikzcd}
% 	{\arr A} && A \\
% 	& {\arr \arr A} & {\arr A} \\
% 	{\arr A \times \arr A} & {\arr (A \times A)} & {A \times A}
% 	\arrow["\src", from=1-1, to=1-3]
% 	\arrow["{l \circ b}"{description}, dashed, from=1-1, to=2-2]
% 	\arrow["{(\rfl \circ \src, \id)}"', from=1-1, to=3-1]
% 	\arrow["\rfl", from=1-3, to=2-3]
% 	\arrow["\src", from=2-2, to=2-3]
% 	\arrow["{\arr (\src,\trg)}"', from=2-2, to=3-2]
% 	\arrow["{(\src,\trg)}", from=2-3, to=3-3]
% 	\arrow["\iso"{description}, draw=none, from=3-2, to=3-1]
% 	\arrow["\src"', from=3-2, to=3-3]
% \end{tikzcd}\]
  \[ 
    \begin{tikzcd}
        {\arr A} & P & {\arr \arr A} & {\arr \arr A}
        \arrow["b", from=1-1, to=1-2]
        \arrow["\chi"', bend right, dashed, from=1-1, to=1-4]
        \arrow["l", from=1-2, to=1-3]
        \arrow["\sym", from=1-3, to=1-4]
    \end{tikzcd}
  \]
\end{proof}

\section{Identity elimination}

\begin{theorem}\label{thm:id-elim}
  Suppose $(\CC,\RR)$ is a $\pi$-preclan.
  If $\uu : \Ulow \to \U$ is an $\RR$-algebraic universe (\labelcref{def:algebraic-universe}) that
  admits $\Path$-types (\labelcref{def:mltt-alg-path-types-2})
  and a normal Hurewicz structure (\labelcref{def:hurewicz-defs}),
  then the universe also admits $\Id$-types (\labelcref{def:algebraic-id-types}).
\end{theorem}
\begin{proof}
  Identity formation and introduction (\Cref{def:algebraic-id-form-intro})
  are provided by the following diagram.
  \[\begin{tikzcd}
    \Ulow & {{\arr[\U] \Ulow}} & \Ulow \\
    & {\Ulow \times_\U \Ulow} & \U
    \arrow["\rfl", from=1-1, to=1-2]
    \arrow["\Delta"', from=1-1, to=2-2]
    \arrow["\path", from=1-2, to=1-3]
    \arrow[from=1-2, to=2-2]
    \arrow["\lrcorner"{anchor=center, pos=0.125}, draw=none, from=1-2, to=2-3]
    \arrow[from=1-3, to=2-3]
    \arrow["\Path"', from=2-2, to=2-3]
  \end{tikzcd}\]
  Note that $I = {\arr[\U] \Ulow}$ is provided algebraically by the path functor.
  
  Let $\Gamma \in \CC$;
  we again drop subscripts for $\arr[\Gamma]$ and always work over $\Gamma$.
  Let $A \to \Gamma$ and $d : C \to \arr A$ be
  two small fibrations,
  and $r : A \to C$ be a map satisfying $d \circ r = \rfl$.
  By \Cref{prop:id-elim-equiv},
  it suffices construct a lift $j : \arr A \to C$
  such that $d \circ j = \id$ and $j \circ \rfl = r$.
  \[ \begin{tikzcd}
	{{A}} & C \\
	{\arr{A}} & {\arr{A}}
	\arrow["r",  from=1-1, to=1-2]
	\arrow["{{{{\rfl}}}}"', from=1-1, to=2-1]
	\arrow[from=1-2, to=2-2]
	\arrow["{{j}}"{description}, dashed, from=2-1, to=1-2]
	\arrow[equals, from=2-1, to=2-2]
  \end{tikzcd} \]
  By \Cref{cor:small-fibration-hurewicz}, $d : C \to \arr A$ is Hurewicz over $\Gamma$.
  To construct $j$, let us consider the following lifting problem in $C$
  for some path $p : x \rightsquigarrow y$ in $A$.
  \[ \begin{tikzcd}
	r & {j_p} & { C} \\
	{\rfl x} & p & { \arr A}
	\arrow[dashed, from=1-1, to=1-2]
	\arrow[maps to, from=1-1, to=2-1]
	\arrow["d", from=1-3, to=2-3]
	\arrow["{\chi_p}"', squiggly, from=2-1, to=2-2]
    \end{tikzcd}\]
  More formally, the lifting problem is captured by the following dotted map.
  \[\begin{tikzcd}
        {\arr A} && A \\
        & P & C \\
        & {\arr \arr A} & {\arr A}
        \arrow["\src", from=1-1, to=1-3]
        \arrow["{(\chi,r \circ \src)}"{description}, dashed, from=1-1, to=2-2]
        \arrow["\chi"', bend right, from=1-1, to=3-2]
        \arrow["r", from=1-3, to=2-3]
        \arrow[from=2-2, to=2-3]
        \arrow["\pi_1", from=2-2, to=3-2]
        \arrow["\lrcorner"{anchor=center, pos=0.125}, draw=none, from=2-2, to=3-3]
        \arrow["d", from=2-3, to=3-3]
        \arrow["\src"', from=3-2, to=3-3]
    \end{tikzcd}\]
    By \Cref{prop:hurewicz-equiv-defs},
    $\arr C$ comes with a locally Hurewicz lifting operation $l : P \to \arr C$.
    We obtain $j$ by composing
    $(\chi,r \circ \src) : \arr A \to P$
    with $l : P \to \arr C$
    and the target map $\trg : \arr C \to C$.
    \[ \begin{tikzcd}
        {\arr A} & P & {\arr C} & C
        \arrow["{(\chi,r \circ \src)}", from=1-1, to=1-2]
        \arrow["j"', bend right, dashed, from=1-1, to=1-4]
        \arrow["l", from=1-2, to=1-3]
        \arrow["\trg", from=1-3, to=1-4]
    \end{tikzcd} \]
    The lower triangle involving $j$ commutes:
    \begin{align*}
      & d \circ j \\
      = \; & d \circ {\trg} \circ l \circ (\chi, r \circ \src) \\
      = \; & {\trg} \circ \arr (d) \circ l \circ (\chi, r \circ \src) \\
      = \; & {\trg} \circ \pi_1 \circ (\chi, r \circ \src) \\
      = \; & {\trg} \circ \chi \\
      = \; & \id \\
    \end{align*}
    The upper triangle commutes due to normality of $\chi$ and $l$.
    \begin{align*}
      & j \circ \rfl \\
      = \; & {\trg} \circ l \circ (\chi, r \circ \src) \circ \rfl \\
      = \; & {\trg} \circ l \circ (\chi\circ \rfl, r \circ \src\circ \rfl)  \\
      = \; & {\trg} \circ l \circ (\rfl \circ \rfl, r) & (\chi \text{ normal}) \\
      = \; & {\trg} \circ l \circ h \circ \rfl \circ r   \\
      = \; & {\trg} \circ \rfl \circ r & (l \text{ normal}) \\
      = \; & r  \\
    \end{align*}
\end{proof}

\section{Examples} \label{sec:examples}

\begin{example}\label{example:sets}
  Take $\RR$ to be all maps in a lex category $\CC$
  and the cylinder functor $\cyl : \CC \to \CC$
  to be the identity functor on $\CC$, 
  then the path functors $\arr[X] : \RR(X) \to \RR(X)$
  will also be identities
  and all maps are trivially normal Hurewicz.
  It follows that the pullback square appearing in \Cref{def:mltt-alg-path-types-2}
  is the same pullback square defining \emph{extensional} Identity types (cf.~\cite{awodey2025}):
  \[ \begin{tikzcd}
        \Ulow & \Ulow \\
        {\Ulow \times_\U \Ulow} & \U
        \arrow["\path", from=1-1, to=1-2]
        \arrow["{{(\src,\trg) = \Delta}}"', from=1-1, to=2-1]
        \arrow["\lrcorner"{anchor=center, pos=0.125}, draw=none, from=1-1, to=2-2]
        \arrow[from=1-2, to=2-2]
        \arrow["\Path"', from=2-1, to=2-2]
    \end{tikzcd} \]

  In particular when $\CC = \Psh(\DD)$ is a category of presheaves,
  $\RR$ is the class of all maps,
  $\uu : \Ulow \to U$ is a Hofmann-Streicher universe \cite{hofmann1997},
  such a pullback square can always be constructed (see \cite[Section {7.1}]{awodey2026}).
\end{example}

As shown in \Cref{ex:exponentiable-bipointed},
a bipointed exponentiable object
$\delta_0, \delta_1 : 1 \rightrightarrows I$
naturally gives rise to a cylinder structure on $\CC$.

\begin{example}\label{example:groupoids}
  Take $\CC$ to be the category of groupoids
  and $I = \{\star \iso \star\}$ to be the walking isomorphism.
  Then the normal Hurewicz condition is equivalent to being a normal isofibration.
  Taking $\RR$ to be the class of isofibrations,
  which is a $\pi$-clan,
  it is clear that exponentiation by $I_X$ provides a partial right adjoint
  $\RR(X) \to \RR(X)$ -- it also extends to an actual right adjoint.
  Taking $\uu : \Ulow \to \U$ to be the universal small isofibration in groupoids
  (see \Cref{sec:universe-in-groupoid-model}),
  it is easy to check the equivalent conditions of \Cref{prop:universal-pathobject}.
  In fact, pathobject projection for a small isofibration
  is a small discrete isofibration.
\end{example}

\begin{example} \label{example:cSet}
  If we take $\CC = \cSet$ to be the category of
  Cartesian cubical sets, $I = y(\Box_1)$,
  and $\RR$ to be the class of unbiased cubical fibrations
  \cite[Definition 4.6]{awodey2026cartesian},
  then $(\CC,\RR)$ is only a $\pi$-preclan.
  To see that $(\CC,\RR)$ fails to be a $\pi$-clan,
  note that the interval $I$ is not fibrant.
  However, for each $\Gamma$, taking $\arr[\Gamma] A := A^{I_\Gamma}$
  still defines an endofunctor $\arr[\Gamma] : \RR(\Gamma) \to \RR(\Gamma)$,
  where $I_\Gamma$ is the pullback of $I$ to the slice $\CC / \Gamma$.
  This is because although $I \to 1$ is not a fibration,
  it is $\RR$-exponentiable.
  We can take $\uu : \Ulow \to \U$
  to be the classifier for small unbiased fibrations
  \cite[Proposition 7.17]{awodey2026cartesian};
  checking that $\U$ itself is fibrant \cite[Proposition 9.4]{awodey2026cartesian}.
  In this case, all unbiased cubical fibrations are Hurewicz,
  and a straightforward computation shows that the universal pathobject projection
  is a small unbiased fibration.
\end{example}

\begin{example} \label{example:sSet}
  Take $\CC = \mathsf{sSet}$ to be the category of
  simplicial sets and $\RR$ to be the class of Kan fibrations.
  This again forms a $\pi$-preclan (\Cref{ex:sset-kan}).
  Exponentiation by $I = [1]$ (and its pullbacks)
  again defines an endofunctor $\arr[1] : \RR(1) \to \RR(1)$.
  Then we can take $\uu : \Ulow \to \U$
  to be the universe from \cite[Definition 2.1.9]{kapulkin2021},
  checking that $\U$ is fibrant \cite[Theorem 2.2.1]{kapulkin2021}.
  Here, it is also the case that all Kan fibrations are Hurewicz.
  A proof that the universal pathobject projection
  is a small unbiased fibration can be adapted from \cite[Proposition 2.3.3]{kapulkin2021}.
\end{example}

\section{A path types modality}
There is a modal character to path types, which we make precise now
by defining a modal theory and showing that our examples are models
of this modal theory.
This suggests a modal type theory that
has characteristics of both (plain) Homotopy Type Theory
and (various versions of) cubical type theory.
Here we will suggest what a model of such a modal type theory should satisfy,
leaving the design of a modal syntax to future work.
% So far,
% we have presented path types in a style that is reminiscent of modal type theory.
% Here, we make this precise by providing
% a modal type theory and an additional condition
% that makes our setting a model of it.

The modal theory is given as follows,
in the style of \cite{licata2016,licata2017,gratzer2020}.
Consider the 2-category $\MM$ (the \emph{mode theory}) with a single 0-cell (\emph{mode}) $\star$,
a single non-trivial 1-cell (\emph{modality}) $\arr : \star \to \star$,
and two-cells $\src, \trg : \arr \to \id$,
$\rfl : \id \to \arr$ and $\sym : \arr \circ \arr \to \arr \circ \arr$.

Let us again assume that
$(\CC,\RR)$ is a $\pi$-preclan with a terminal object,
a cylinder structure and path functors (\Cref{def:path-functor}).
It is straightforward to define a strict 2-functor
$\llbracket - \rrbracket : \MM^\mathsf{coop} \to \tcat$
(cf. \cite[Definition 5.1]{gratzer2020}),
using the cylinder structure.
The mode $\star$ is interpreted as the category $\CC$,
the modality $\arr$ is interpreted as the cylinder functor $\cyl : \CC \to \CC$
and the two-cells are interpreted as the obvious natural transformations.

\begin{definition}\label{def:modal-model}
  A modal universe with respect to the context structure
  $\llbracket - \rrbracket : \MM^\mathsf{coop} \to \tcat$
  consists of an $\RR$-algebraic universe $\uu : \Ulow \to \U$,
  maps $\Mod : \arr[1] \U \to \U$ and $\mod : \arr[1] \Ulow \to \Ulow$,
  such that
  \[\begin{tikzcd}
      {\arr[1] \Ulow} & \Ulow \\
      {\arr[1] \U} & {  \U}
      \arrow["\mod", from=1-1, to=1-2]
      \arrow["{\arr[1]\uu}"', from=1-1, to=2-1]
      \arrow["\lrcorner"{anchor=center, pos=0.125}, draw=none, from=1-1, to=2-2]
      \arrow["\uu", from=1-2, to=2-2]
      \arrow["\Mod"', from=2-1, to=2-2]
  \end{tikzcd}\]
  is a pullback square.
\end{definition}
Given an interval object $I$,
the map $\Mod : \arr[1] \U \to \U$ should be viewed as
a kind of $\Pi$-type indexed by $I$,
taking a type $A : I \to \U$ to the type $\Pi_I A : \U$.

Since we are working internally,
the above definition of a modal universe is the internal analogue of a
``modal natural model'' (\cite{gratzer2020}) with a ``dependent right adjoint''
(\cite{birkedal2020}).
By taking the image of $\uu : \Ulow \to \U$ under the Yoneda embedding,
we obtain a model of modal type theory with 
a dependent right adjoint provided by the partial right adjoints
$\arr[\Gamma] : \RR(\Gamma) \to \RR(\Gamma)$
for each $\cyl[\Gamma] : \CC / \Gamma \to \CC / \Gamma$.
One way to state a precise result is using CwDRAs
(categories with dependent right adjoints),
but the same could be done with modal natural models.

\begin{proposition}
  Suppose $(\CC,\RR)$ is a $\pi$-preclan with a terminal object,
  a cylinder structure and path functors (\Cref{def:path-functor}).
  A modal universe provides a model of the modal type theory $\MM$
  in the sense of being a CwDRA (\cite[Definition 2]{birkedal2020})
  where the presheaf of types is representable.
\end{proposition}

% Cylinders on contexts $\cyl : \CC \to \CC$
% are the functors associated with the modality
% $\arr$ and paths on types $\arr[\Gamma] : \RR(\Gamma) \to \RR(\Gamma)$
% are their so-called ``{dependent right adjoints}''.
% As usual,
% each two-cell between modalities induces a natural transformation
% between the associated functors.

If the universe $\uu : \Ulow \to \U$ has classifying maps
(\Cref{def:alg-universe-classifying-map}),
then we can also rephrase the pullback condition in \Cref{def:modal-model} as saying
that the path functor $\arr[1]$ preserves small fibrations
(pullbacks of the universe $\uu$).
\[ \arr[1] : (\RR(1),\RR_\U) \to (\RR(1),\RR_\U) \]
(If $A$ and $B$ are in $\RR(1)$ and $f : A \to B$ is a small fibration then
so is $\arr[1] f : \arr[1] A \to \arr[1] B$.)
This is useful because in practice it is often easier
to check whether the path functor preserves small fibrations,
as opposed to directly constructing the maps $\Mod$ and $\mod$.
To demonstrate this,
let us revisit the previous examples.

\begin{example}
  % The examples from before,
  % using an interval object $I$,
  % all admit modal universes,
  % and are therefore models of the model type theory
  % given by $\MM$.
  In $\Set$, taking $I = 1$ as the interval,
  exponentiation by $1$ is the identity,
  and so path functors preserve small fibrations.

  In $\Grpd$ with $I$ as the walking isomorphism,
  exponentiation by $I$ preserves small isofibrations because
  $I \to 1$ is itself a small isofibration,
  and small isofibrations form a $\pi$-preclan.
  One way to see this is by using \Cref{lem:pushforward-preserves-maps},
  with $\SS$ as small isofibrations and $\RR$ as the
  $\pi$-clan of isofibrations.

  In $\mathsf{cSet}$ with $I = y(\Box_1)$,
  although $I \to 1$ is not a fibration,
  it is $\RR$-exponentiable (see \Cref{ex:cSet-unbiased-fibrations}),
  where $\RR$ is the class of unbiased fibrations.
  It follows that exponentiating by $I$ preserves small unbiased fibrations:
  again we can use \Cref{lem:pushforward-preserves-maps},
  with $\SS$ as small unbiased fibrations and $\RR$ as the
  $\pi$-preclan of unbiased fibrations.
\end{example}

% \begin{remark}
%   We have only used the $\pi$-clan for making polynomial functors,
%   which in turn are used to construct the generic case for identity elimination.
%   As was demonstrated in \Cref{sec:exponentiability} and \Cref{sec:polynomials},
%   a $\pi$-clan is not the setting in which polynomials can be defined.
%   A similar exponentiability condition was omitted in \cite{awodey2026},
%   where coherence of identity elimination was not proven;
%   one way to fix this issue is to require that the universe is exponentiable.
%   % As noted in \Cref{a},
%   % $\pi$-clans are not the only way one can obtain polynomial functors

%   % If such a condition were added,
%   % then the axioms presented here are more general than those in \cite{awodey2026}.
% \end{remark}